\section{全链条、核心池与需求锚点}

本案例建立了43个节点的智能汽车电子全链池：上游包括SoC、存储、显示、感知、连接和基础软件；中游包括德赛西威与国内外Tier-1竞争者；下游包括自主、合资和海外OEM。节点池刻意保留海外、私营、低纯度和不可得节点，以免把“没有证据的关联”从图上删除后误变成确定关系。

核心估值池只有德赛西威。华阳集团、均胜电子、科博达、经纬恒润、Bosch、Continental、Aptiv和Visteon是竞争参照；Qualcomm、NVIDIA及国产平台候选是上游技术/供给锚点；理想、小米、吉利、奇瑞、长城、广汽丰田等OEM是客户和需求锚点。卫星和需求节点都不进入本案的收入或估值计算。

\begin{exhibitbox}[核心、卫星与需求锚点]
\small
\begin{tabularx}{\textwidth}{L{2.4cm}Y Y Y}
\toprule
类别 & 代表节点 & 已知关系 & 估值处理 \\
\midrule
核心估值池 & 德赛西威 & 已有座舱、智驾和网联分部收入；部分客户/平台量产事实 & 仅按汇总分部历史和保守情景估值 \\
卫星竞争池 & 华阳、均胜、科博达、经纬恒润及全球Tier-1 & 同类产品/客户竞争逻辑，具体项目重叠未核验 & 竞争参照，不给客户链信用 \\
上游平台 & Qualcomm、NVIDIA、国产SoC候选、显示/感知供应商候选 & 特定合作或历史项目可验证，其余采购关系未披露 & 供给/替代风险，不倒推成本或收入 \\
OEM需求锚点 & 理想、小米、吉利、奇瑞、广汽等 & 年报列示的配套、量产或订单关系 & 需求/项目证据，不证明客户收入份额 \\
\bottomrule
\end{tabularx}
\sourcenote{全链节点池、公司年报、合作方官方材料。}
\end{exhibitbox}

这个分类的关键是防止“下游强需求”自动变成“上游确认收入”。OEM销量、平台发布和品牌热度只能提升验证优先级；它们无法回答德赛西威到底供应哪个车型、何时SOP、单车价格多少、毛利多少、回款多快。相反，公司已确认的分部收入与毛利才是当前可用的盈利基础。

\section{客户、平台与订单的证据质量}

第四代座舱在理想、小米、吉利、奇瑞、广汽等客户配套是较强的客户/产品量产事实。智能驾驶域控在小米、理想、长城、小鹏、吉利、广汽丰田、东风日产、奇瑞、广汽本田等大规模量产的年报表述，支持公司具有多客户系统交付能力。Qualcomm G10PH合作和NVIDIA/小鹏P7历史协作则说明公司具有多平台适配和历史项目能力。

然而，每一类证据都有边界。中央计算订单没有实名客户、金额和收入；AR-HUD/区域控制器订单缺SOP/收入；Qualcomm合作没有OEM量产收入；NVIDIA的历史项目没有当前持续供货信息。正确的模型处理是把这些字段留为not disclosed，设置升级路径，而非从客户名称推导收入。

\begin{exhibitbox}[客户链审计：什么可以进入基准]
\small
\begin{tabularx}{\textwidth}{L{2.6cm}Y Y Y}
\toprule
关系/产品 & 可采用表述 & 缺失字段 & 基准估值处理 \\
\midrule
第四代智能座舱 & 已在多个具名客户配套 & 车型、ASP、客户收入、项目毛利 & 已确认座舱汇总收入/毛利可用；不拆客户 \\
智驾域控 & 多OEM大规模量产 & 客户份额、SoC/BOM、项目利润与回款 & 已确认智驾收入/毛利可用；不加客户驱动EPS \\
中央计算/第五代座舱 & 匿名头部客户订单/开发中 & 客户、SOP、金额、收入、毛利 & 零增量估值信用 \\
AR-HUD/区域控制器 & 有订单，部分即将量产 & 车型、量产时点、ASP、收入 & 零增量估值信用 \\
平台合作 & 特定合作或历史项目可验证 & 当前采购、独供、未来车型收入 & 风险/能力证据，不作收入输入 \\
\bottomrule
\end{tabularx}
\sourcenote{客户链审计、公司年报、Qualcomm及NVIDIA官方资料。}
\end{exhibitbox}

\section{价值量、毛利与现金转换}

价值链的核心不是“组件有没有价值”，而是价值能否在Tier-1处转化为收入、毛利和现金。2025年，座舱收入CNY20.585bn、毛利率18.83\%；智驾收入CNY9.700bn、毛利率16.36\%；汽车电子原材料CNY24.183bn，占营业成本91.78\%。这说明BOM、芯片/显示/感知成本、项目报价和产品结构都可能显著影响利润。

公司披露工厂投产与建设，能够说明区域交付能力；没有设计产能、利用率、良率、项目订单和价格，就不能用工厂面积或建设进度估算收入。类似地，库存增长和订单备货可以解释经营过程，不能跳过后续车型销售、验收和回款的风险。

\begin{riskbox}[产业链估值底线]
客户、订单、ASP、利用率和项目毛利任一项缺失时，不把新增平台的潜在收入计入目标价。德赛西威仍可基于已确认分部进行汇总估值；但由客户链推导的增量模型维持观察名单状态。
\end{riskbox}
